\documentclass[12pt]{article}
\usepackage{amsmath}
\usepackage{amssymb}
\usepackage{geometry}
\geometry{margin=1in}

\title{Lecture Scribe}
\author{Course: CSE400 – Fundamentals of Probability in Computing}
\date{Lecture: L17 \\ Topic: Joint Random Variables and Joint Distributions \\ Instructor: Dhaval Patel, PhD \\ Date: March 10, 2026}

\begin{document}

\maketitle
\noindent\rule{\linewidth}{0.4pt}
\section{Definitions and Notation}

\textbf{Joint Random Variables}

In many applications, two or more random variables are considered simultaneously.\\
Examples of random variables considered jointly:

\begin{itemize}
\item $X =$ Vehicle Speed
\item $Y =$ Traffic Density
\item $X =$ Signal-to-Noise Ratio (SNR)
\item $Y =$ Data Throughput
\item $X =$ Rainfall
\item $Y =$ Temperature
\item $X =$ Wind Speed
\item $Y =$ Delivery Time
\item $X, Y =$ outcomes of rolling two dice
\end{itemize}
\noindent\rule{\linewidth}{0.4pt}
\textbf{Joint Probability Mass Function (Discrete Case)}

For discrete random variables $X$ and $Y$, the joint PMF represents the probability that $X$ and $Y$ simultaneously take specific values.\\
{Notation:}

\[
p_{X,Y}(x,y) = P(X = x, Y = y)
\]
\newpage
{Representation:}

\begin{itemize}
\item Joint PMF Table Representation
\end{itemize}
\noindent\rule{\linewidth}{0.4pt}

\textbf{Joint Probability Density Function (Continuous Case)}

For continuous random variables $X$ and $Y$, the joint distribution is described using a \textbf{joint probability density function}.\\
Notation:
\[
f_{X,Y}(x,y)
\]
\noindent\rule{\linewidth}{0.4pt}
\textbf{Joint Cumulative Distribution Function}

The joint CDF of two random variables $X$ and $Y$ is defined as

\[
F_{X,Y}(x,y) = P(X \le x, Y \le y)
\]
\noindent\rule{\linewidth}{0.4pt}
\section{Assumptions}

Situations involving multiple random variables may include relationships or correlations between the variables.

Examples provided in the lecture include:

\begin{itemize}
\item Vehicle speed and traffic density
\item Signal-to-noise ratio and throughput
\item Rainfall and temperature
\item Wind speed and delivery time
\end{itemize}
\noindent\rule{\linewidth}{0.4pt}
\section{Theorem / Results}

\textbf{Joint PMF Representation}\\

A joint PMF may be represented in tabular form where:

\begin{itemize}
\item rows correspond to values of $X$
\item columns correspond to values of $Y$
\item table entries represent $P(X=x, Y=y)$
\end{itemize}
\noindent\rule{\linewidth}{0.4pt}
\textbf{Joint CDF Definition}\\
$F(x,y) = P(X \le x, Y \le y)$\\
\noindent\rule{\linewidth}{0.4pt}
\textbf{Joint CDF -- Continuous Case}

The joint CDF for continuous random variables is obtained from the joint PDF.\\
\noindent\rule{\linewidth}{0.4pt}
\textbf{Joint CDF -- Discrete Case}

For discrete random variables, the joint CDF is obtained from the joint PMF.\\
\noindent\rule{\linewidth}{0.4pt}
\textbf{Properties of the Joint CDF}

The lecture introduces properties associated with the joint cumulative distribution function.\\
\noindent\rule{\linewidth}{0.4pt}
\section{Proof / Derivation}

No explicit proofs or derivations are formally written in the provided lecture slides.\\
\noindent\rule{\linewidth}{0.4pt}
\section{Worked Examples}

\textbf{Example 1 -- Smart Transportation}

Random variables:

\begin{itemize}
\item $X =$ Vehicle Speed
\item $Y =$ Traffic Density
\end{itemize}

Question:

Can speed be high when the traffic density is also high?

This illustrates correlated random variables.\\
\noindent\rule{\linewidth}{0.4pt}
\textbf{Example 2 -- Wireless Network Performance}

Random variables:

\begin{itemize}
\item $X =$ Signal-to-Noise Ratio (SNR)
\item $Y =$ Data Throughput
\end{itemize}

Question:

If SNR is high, will throughput always be high?\\
\noindent\rule{\linewidth}{0.4pt}
\newpage
\textbf{Example 3 -- Weather Prediction}

Random variables:

\begin{itemize}
\item $X =$ Rainfall
\item $Y =$ Temperature
\end{itemize}

Application: weather forecasting.\\
\noindent\rule{\linewidth}{0.4pt}

\textbf{Example 4 -- Drone Delivery}

Random variables:

\begin{itemize}
\item $X =$ Wind Speed
\item $Y =$ Delivery Time
\end{itemize}

Application: logistics and UAV systems.\\
\noindent\rule{\linewidth}{0.4pt}

\textbf{Example 5 -- Rolling Two Dice}

Two random variables:

\begin{itemize}
\item $X =$ outcome of first die
\item $Y =$ outcome of second die
\end{itemize}

\textbf{Example conditions:}

\begin{itemize}
\item $X + Y = 7$
\item $X > Y$
\end{itemize}

Joint PMF is required for such events.\\
\noindent\rule{\linewidth}{0.4pt}


\section{Conclusion}

The lecture introduces:

\begin{itemize}
\item Joint random variables
\item Joint PMF for discrete variables
\item Joint PDF for continuous variables
\item Joint Cumulative Distribution Function
\item Applications involving multiple random variables
\end{itemize}
\noindent\rule{\linewidth}{0.4pt}
\begin{center}
\textbf{End of Lecture Scribe}
\end{center}
\end{document}